\section*{الفصل \ref{ch:g12:early-earth} --- الأرض الفتية وغلافها الجوي}

\begin{solution}{exo:g12:early-earth:1}
(أ) \omterm{def:g12:early-earth:primitive-atmosphere}{انطلاق الغازات}؛ (ب) \omterm{def:g12:early-earth:stromatolite}{الستروماتوليت}؛ (ج) \omterm{def:g12:early-earth:bif}{تكوين الحديد الشريطي}؛ (د) \omterm{def:g12:early-earth:red-beds}{الطبقات
الحمراء}؛ (ه) \omterm{def:g12:early-earth:goe}{حدث الأكسدة العظيم}.
\end{solution}

\begin{solution}{exo:g12:early-earth:2}
حبيبات الزركون التي تُظهر الماء السائل (4.4 مليار سنة)؛ فأولى \omterm{def:g12:early-earth:stromatolite}{الستروماتوليتات}
في بيلبارا (3.43)؛ \omterm{def:g12:early-earth:bif}{فتكوينات الحديد الشريطية} في هامرسلي (أكثر من
2.5)؛ فنهاية \omterm{def:g12:early-earth:goe}{حدث الأكسدة العظيم} (2.06)؛ فنهاية معظم \omterm{def:g12:early-earth:bif}{تكوينات الحديد
الشريطية} (1.8).
\end{solution}

\begin{solution}{exo:g12:early-earth:3}
عند $t = 0$، $L = L_\odot/(1 + 0.4) = L_\odot/1.4 \approx 0.714\,L_\odot$:
أي أخفت بنحو \qty{29}{\percent}، وأقل سطوعًا بقرابة الثلث مما هي اليوم.
\end{solution}

\begin{solution}{exo:g12:early-earth:4}
أن الهواء كان حينئذ لا يكاد يحوي أكسجينًا. فالبيريت يدمره الأكسجين؛
ولم يكن لحبيباته أن تتدحرج مع الأنهار وتترسب سليمة إلا إذا خلا منه
الهواء والماء عند السطح.
\end{solution}

\begin{solution}{exo:g12:early-earth:5}
بخار الماء، الذي تكاثف فصار المحيط، وثاني أكسيد الكربون، الذي
بقي في الهواء بكميات أكبر بكثير مما اليوم ودفّأ الأرض
بأثر الدفيئة.
\end{solution}

\begin{solution}{exo:g12:early-earth:6}
قبل 3 مليارات سنة: $1 - t/t_\odot = 3/4.567 \approx 0.657$، و$L \approx
1/1.263 \approx 0.792\,L_\odot$، و$T \approx 255 \times 0.792^{1/4} \approx
\qty{241}{K}$. وقبل 2.45 مليار سنة: $1 - t/t_\odot \approx 0.536$، و$L
\approx 1/1.215 \approx 0.823\,L_\odot$، و$T \approx 255 \times 0.823^{1/4}
\approx \qty{243}{K}$.
\end{solution}

\begin{solution}{exo:g12:early-earth:7}
\qty{1}{kg} من الحديد يساوي $1000/55.85 \approx 17.9$ مول؛ وجزيء واحد من
الأكسجين يأخذ أربعة أيونات من الحديد، إذن $17.9/4 \approx 4.48$ مول من
\ce{O2}، أي $4.48 \times 32.00 \approx \qty{143}{g}$ من الأكسجين.
\end{solution}

\begin{solution}{exo:g12:early-earth:8}
(1) $2.46 - 2.06 = 0.40$ مليار سنة، أي 400 مليون سنة. (2) من 4.54
إلى 2.46 مليار سنة مضت: $2.08/4.54 \approx \qty{46}{\percent}$ من
تاريخ الأرض. (3) نحو $3.43 - 2.46 \approx 1$ مليار سنة: فالأكسجين
الذي صنعته الميكروبات ظل يُستهلك زمنًا طويلًا قبل أن يستطيع البقاء في
الهواء.
\end{solution}

\begin{solution}{exo:g12:early-earth:9}
\omterm{def:g12:early-earth:bif}{تكوين الحديد الشريطي} يدل على بحر عميق خالٍ من الأكسجين، كان يستطيع حمل
الحديد الذائب، تحت بحر ضحل مشمس فيه من الأكسجين ما يكفي لترسيبه.
أما الطبقة الحمراء، المترسبة على اليابسة، فتدل على أكسجين في الهواء وفي الماء
عند السطح، يكفي لإصداء حديد حبيباتها.
\end{solution}

\begin{solution}{exo:g12:early-earth:10}
في بركة هاملين، خليج القرش، غرب أستراليا، حيث الماء أملح بمرتين
من ماء البحر المعتاد. يحتجز بساط من الميكروبات، أساسًا البكتيريا الزرقاء،
حبيبات الرسوبيات ويربطها ويجعل الكربونات تنمو؛ ثم ينمو بساط جديد فوقه،
وهكذا، طبقة فوق طبقة.
\end{solution}

\begin{solution}{exo:g12:early-earth:11}
(1) $21 \times 10^{-5} = \qty{0.00021}{\percent}$. (2) $21/10 =
\qty{2.1}{\percent}$.
\end{solution}

\begin{solution}{exo:g12:early-earth:12}
تبقى \omterm{def:g10:habitable-planet:equilibrium-temperature}{درجة حرارة التوازن} دون تغيير إذا ساوى $L(1 - A')$ القيمة $1 - A$:
$1 - A' = 0.706/0.741 \approx 0.953$، إذن $A' \approx 0.047$. فكان على الأرض
أن تعكس أقل من \qty{5}{\percent} من ضوء الشمس، كأنها تقريبًا
جسم أسود، وهذا ما لا يقدر عليه كوكب فيه محيط وسحب وبعض الجليد.
والمخرج الآخر أثر دفيئة أقوى.
\end{solution}

\begin{solution}{exo:g12:early-earth:13}
(1) $273 - 237 = \qty{36}{K}$ لبلوغ التجمد؛ و$288 - 237 = \qty{51}{K}$
لبلوغ متوسط اليوم. (2) أثر الدفيئة اليوم يضيف $288 - 255 =
\qty{33}{K}$: فالأرض الفتية كانت تحتاج إلى أثر أقوى. (3) \qty{0.1}{bar}
عُشر ضغط الغلاف الجوي الحالي كله، لثاني أكسيد الكربون وحده.
\end{solution}

\begin{solution}{exo:g12:early-earth:14}
(1) قبل 2450 مليون سنة مضت كانت البصمة موجودة، فكان الهواء يحوي أقل
من $10^{-5}$ من أكسجين اليوم؛ وبعد 2090 مليون سنة غابت، فكان
الهواء يحوي أكثر. (2) نعم: \omterm{def:g12:early-earth:goe}{فحدث الأكسدة العظيم} يمتد من نحو 2460 إلى
2060 مليون سنة مضت، وهي الفترة نفسها.
\end{solution}

\begin{solution}{exo:g12:early-earth:15}
(1) نحو $4567 - 60 = 4507$ مليون سنة مضت. (2) على الأكثر $4507 - 4404 =
103$ مليون سنة بعد الاصطدام: فقد كانت للأرض قشرة وماء في غضون
نحو مئة مليون سنة من الاصطدام.
\end{solution}

\begin{solution}{pb:g12:early-earth:1}
\textbf{1.} $4.567 - 4.0 = 0.567$ مليار سنة.

\textbf{2.} $1 - 0.567/4.567 \approx 1 - 0.124 = 0.876$.

\textbf{3.} $L = 1/(1 + 0.4 \times 0.876) = 1/1.350 \approx 0.741$ من
لمعانها اليوم.

\textbf{4.} $0.741 \times 1361 \approx \qty{1008}{W/m^2}$.

\textbf{5.} $1 - 0.741 \approx 0.26$: أي أخفت بنسبة \qty{26}{\percent}.

\textbf{6.} $T_{\mathrm{eq}} = \left(S(1-A)/(4\sigma)\right)^{1/4}$.

\textbf{7.} مع ثبات $A$ و$\sigma$، تكون $T_{\mathrm{eq}}$ ثابتًا مضروبًا في
$S^{1/4}$، إذن $T'/T = (S'/S)^{1/4}$.

\textbf{8.} $0.741^{1/4} = \sqrt{\sqrt{0.741}} \approx \sqrt{0.861} \approx
0.928$.

\textbf{9.} $255 \times 0.928 \approx \qty{237}{K}$.

\textbf{10.} $237 - 273 \approx \qty{-36}{\celsius}$.

\textbf{11.} أنه كان عند السطح ماء سائل، أي محيط، قبل 4.4
مليار سنة.

\textbf{12.} في ماء ضحل، كما تنمو \omterm{def:g12:early-earth:stromatolite}{الستروماتوليتات} الحية اليوم.

\textbf{13.} \omterm{def:g12:early-earth:bif}{تكوينات الحديد الشريطية} في سلاسل هامرسلي، التي ترسبت
على قاع بحر.

\textbf{14.} نعم: فالدرجة \qty{237}{K} أدنى بكثير من \qty{273}{K}؛ ومع الغلاف
الجوي الحالي كان متوسط السطح سيكون متجمدًا.

\textbf{15.} تُظهر الصخور ماءً سائلًا تحت شمس أخفت من أن تبقيه
سائلًا مع الغلاف الجوي الحالي: وهذه \omterm{def:g12:early-earth:faint-young-sun}{مشكلة الشمس الفتية الخافتة}.

\textbf{16.} $288 - 255 = \qty{33}{K}$.

\textbf{17.} على الأقل $273 - 237 = \qty{36}{K}$، أي أكثر مما اليوم.

\textbf{18.} $A' = 1 - 0.706/0.741 \approx 0.047$: كوكب أسود تقريبًا،
وهذا غير معقول.

\textbf{19.} ثاني أكسيد الكربون، الذي كان أوفر بكثير مما اليوم؛ وربما أسهم الميثان،
الذي صنعته الميكروبات الأولى.

\textbf{20.} كانت \omterm{def:g10:habitable-planet:equilibrium-temperature}{درجة حرارة التوازن} للأرض قبل أربعة مليارات سنة
\textbf{نحو 237 كلفن}، أي نحو \qty{36}{\celsius} تحت درجة التجمد.
\end{solution}
